\section{Introduction}
\label{sec:introduction}

Modern language model evaluation infrastructure is dominated by two
complementary modes. Capability benchmarks such as HELM
\citep{liang2022holistic}, MMLU \citep{hendrycks2020measuring},
and BigBench \citep{srivastava2022beyond} answer the question
\emph{``how does model~A compare to model~B on task~$X$?''} They are
indispensable for model selection and capability tracking. Red teaming
and adversarial evaluation \citep{perez2022red, ganguli2022red} answer
\emph{``what can be elicited from this model under pressure?''} They are
indispensable for safety and policy work.

Neither of these modes answers a third question that arises constantly
in deployment: a user has issued a specific request to a specific model
in a specific context, and the model has produced an unexpected
response --- a refusal that seems unjustified, a sycophantic agreement
to a fragile claim, a hedged answer that omits something the user
clearly wanted, or a tone shift mid-conversation. The user wants to
know \emph{why}. Was it the model's own training? A system prompt the
user cannot inspect? An anchor planted earlier in the same
conversation? A keyword filter in the runtime layer? Capability
benchmarks aggregate over thousands of examples and cannot answer this.
Red teaming methodology is adversarial by design and does not fit a
collaborative debugging context.

\paragraph{Robopsychology.}
We introduce \emph{robopsychology}, a behavioral diagnostic toolkit
that targets exactly this third mode of inquiry. The name is borrowed
from Isaac Asimov, whose fictional discipline imagined diagnosing
emergent behavior in machines that follow formal rules without
reprogramming them \citep{asimov1950}. The toolkit's premise is the
same: explanation through structured behavioral interview, not weight
access; diagnosis, not adversarial probing; user-side, not vendor-side.

\paragraph{Contributions.}
This paper makes three contributions:
\begin{enumerate}
    \item A \emph{methodology} (Section~\ref{sec:method}) consisting of
          (i)~a three-way layer split that decomposes any model response
          into model-level, runtime/host, and conversation effects;
          (ii)~explicit \emph{Observed} vs.\ \emph{Inferred} evidence
          labels at the claim level; (iii)~a diagnostic \emph{ratchet}
          --- a fixed-length multi-step probe whose \emph{coherence}
          across steps is the load-bearing signal of whether the model
          is reasoning from a stable internal model or fabricating
          continuity step-by-step.
    \item An \emph{implementation} (Section~\ref{sec:implementation})
          comprising a CLI, a YAML-templated prompt library, and two
          coherence analyzers: a regex baseline and an LLM-judge with a
          four-layer architecture (deterministic pre-filter, judge with
          \texttt{temperature=0} and JSON-mode output, automatic
          regex fallback, exposed severity-graded scoring weights).
    \item Three \emph{reproducible case studies}
          (Section~\ref{sec:validation}) with committed transcripts and
          artifacts. The headline finding (Case~3) is a $0.53$-point
          score delta and a full re-classification from
          \emph{performed} to \emph{genuine} coherence between the
          regex and LLM-judge analyzers on the same 9-step transcript,
          directly validating the LLM-judge approach.
\end{enumerate}

The remainder of the paper is organized as follows. Section~\ref{sec:related}
positions robopsychology against benchmarks, red teaming, alignment evals,
behavioral testing frameworks, and interpretability.
Section~\ref{sec:method} formalizes the method.
Section~\ref{sec:implementation} describes the toolkit.
Section~\ref{sec:validation} presents the three case studies.
Section~\ref{sec:limitations} discusses limitations honestly.
Section~\ref{sec:future} sketches future work.

\todo{Add a final paragraph framing the broader claim: behavioral
diagnostics deserves a seat alongside capability benchmarks and red
teaming, not as a replacement but as a third complementary mode for
practitioners. Keep this paragraph short --- the contributions list
already does most of the framing.}
